\section*{Appendice}

\subsection*{Bottleneck in $\MAXFLOW$ con cammini minimi}
\label{proof:bottleneck_in_max_flow}
Contesto: Riduzione di $\MAXFLOW$ a $\REACHABILITY$ (\ref{riduzione:maxflow_reachability})\\

Se scegliamo sempre il cammino di miglioramento piú corto, 
ogni arco del network residuo puó diventare un bottleneck al massimo $n$ volte.

\begin{proof}
Sia $d_i(v)$ la distanza (numero di archi) da $s$ a $v$ nel grafo residuo dopo l'$i$-esima iterazione. \\
Questa distanza non decresce mai nel corso dell'algoritmo:
$$d_{i+1}(v) \geq d_i(v) \quad \forall v$$
(Intuizione: un'iterazione satura archi in avanti e ne aggiunge all'indietro; 
gli archi indietro aggiunti non possono "accorciare" nessun cammino, 
perché se lo facessero avrebbero già dovuto far parte di un cammino minimo precedente,
altrimenti $d_i$ non sarebbe stato minimo.) \\


Supponiamo che l'arco $(u,v)$ sia bottleneck all'iterazione $i$. \\
Siccome il cammino usato è un cammino minimo,
$$d_i(v) = d_i(u) + 1$$
Essendo bottleneck, $(u,v)$ satura e "scompare" dal grafo residuo (la sua capacità residua diventa 0); 
contemporaneamente compare (o si rafforza) l'arco inverso $(v,u)$.

Perché $(u,v)$ possa tornare disponibile nel grafo residuo — 
e quindi possa ridiventare bottleneck una seconda volta, alla iterazione $j > i$ — 
deve prima essere "riattivato" da un'iterazione che passa per l'arco inverso $(v,u)$. 
In quel momento $(v,u)$ sta sul cammino minimo, quindi:
$$d_j(u) = d_j(v) + 1$$

Applicando le disuguaglianze precedenti, otteniamo:
$$d_j(u) = d_j(v) + 1 \geq d_i(v) + 1 = (d_i(u)+1) + 1 = d_i(u) + 2$$


Ogni volta che l'arco $(u,v)$ (o il suo inverso $(v,u)$) diventa bottleneck, 
la distanza $d(u)$ da $s$ aumenta di almeno 2 rispetto alla volta precedente. \\

Ma $d(u)$ può assumere solo valori tra $0$ e $|V|-1$ 
(al più $|V|$ valori distinti, essendo $u \neq s$ e $u \neq t$ tipicamente tra $1$ e $|V|-2$). 
Siccome $d(u)$ cresce di almeno 2 ad ogni "resurrezione" dell'arco, 
quest'ultimo può diventare bottleneck al più
$O(|V|/2) = O(n)$
\end{proof}
Prova generata da Claude Sonnet 5 (Extra).


\newpage
\subsection*{Raffinamento della relazione tra tempo e spazio}
\label{proof:raffinamento_time_space}
Contesto: limite superiore dello spazio rispetto al tempo (\ref{raffinamento:time_space}). \\
\[ 
\TIME(f(n)) \subseteq \SPACE{ ( \frac{f(n)}{\log(f(n))} ) } 
\]
\begin{proof}
The proof standardizes the computation of a 
$f(n)$ time-bounded Turing machine on input $x$ 
so that the strings are divided in blocks of size $\sqrt{f(|x|)}$. 
The machine is \textit{block-respecting}, that is, 
block boundaries are crossed only in steps 
that are integer multiples of $\sqrt{f(|x|)}$.
\begin{enumerate}

\item (a) Show that any $f(n)$ time-bounded $k$-string Turing machine 
can be made block-respecting with no increase in time complexity.
\subitem The operation of a block-respecting, $k$-string Thring machine $M$ on input $x$ 
can becaptured by a graph $G_M(x)$ with $f(|x|)$ nodes each of indegree and outdegree $O(k)$. \\
The nodes stand for computation segments of length $\sqrt{f(|x|)}$, while edges signify
"information flow." That is, $(B, B')$ is an edge if and only if information from segment
$B$ is necessary for carrying out the computation in segment $B'$.

\item (b) Define precisely the graph $G_M(x)$. Argue that it is acyclic.
\subitem Carrying out the computation of $M$ on $x$ in limited space can n ow be thought of as
"computing" the graph $G_M(x)$, using "registers" of size $\sqrt{f(|x|)}$. The point is to do
it with as few registers as possible, perhaps using recomputation.

\item (c) Define precisely what is a computation of a directed acyclic graph by $R$ registers.
Show that any directed acyclic graph with $n$ nodes and bounded indegree and outdegree 
can be computed with $O( \frac{n}{\log n} )$ registers. 
(This is by far the hardest part of the proof, 
and it involves a sophisticated "divide-and-conquer" technique.)

\item (d) Conclude that $\TIME(f(n)) \subseteq \SPACE{ ( \frac{f(n)}{\log(f(n))} ) }$. 
(Use the (c) above to guess a computation on $G_M(x)$, and apply Savitch's theorem).

\item (e) Show that for one-string machines M, $G_M(X)$ is always planar. What can you
conclude about time and space in one-string machines? (Prove first that planar graphs
can be computed with $O( \sqrt{n} )$ registers.)

\end{enumerate}
\end{proof}
Dal problema 7.4.17 del Papadimitriou, a sua volta tratto dall'articolo \\
J. E. Hopcroft, W. J. Paul, and L. G. Valiant "On time vs.- space and related
problems," Proc. 16th Annual IEEE Symp. on the Foundations of Computer
Sdence, pp. 57-64, 1975